\documentclass[12pt,a4paper]{article}
\usepackage[utf8]{inputenc}
\usepackage{amsmath, amssymb, amsthm}
\usepackage{geometry}
\usepackage{microtype}
\usepackage{parskip}
\usepackage{titlesec}
\usepackage{booktabs}
\usepackage{xcolor}
\usepackage{mdframed}
\usepackage{hyperref}
\usepackage{setspace}

\geometry{margin=2.4cm}

\definecolor{deepblue}{RGB}{10,40,90}
\definecolor{frameblue}{RGB}{200,215,240}

\hypersetup{colorlinks=true, linkcolor=deepblue, urlcolor=deepblue}

\titleformat{\section}{\large\bfseries\color{deepblue}}{}{0em}{}[\vspace{-4pt}\color{deepblue}\hrule\vspace{4pt}]
\titleformat{\subsection}{\normalsize\bfseries\itshape}{}{0em}{}

\newmdenv[
  backgroundcolor=frameblue!25,
  linecolor=deepblue,
  linewidth=0.8pt,
  innerleftmargin=12pt,
  innerrightmargin=12pt,
  innertopmargin=8pt,
  innerbottommargin=8pt,
  roundcorner=3pt
]{highlight}

\newtheorem{remark}{Remark}
\newtheorem{definition}{Definition}

\title{%
  \vspace{-1cm}%
  {\large\scshape Field-Theoretic Political Economy}\\[6pt]
  {\LARGE\bfseries\color{deepblue}The Entropy of Austerity:\\[4pt]
  Modeling the Capital Order via\\[2pt]
  RSVP and Spherepop Dynamics}\\[10pt]
}
\author{Flyxion \\ \small Independent Researcher}
\date{\small April 2026}

\begin{document}
\maketitle
\thispagestyle{empty}

\begin{abstract}
\noindent
This essay synthesizes Clara E.\ Mattei's historical and economic critique of austerity with the Relativistic Scalar-Vector Plenum (RSVP) and Spherepop architectures. Austerity is analyzed not merely as a set of fiscal and monetary policies, but as a rigid, low-entropy configuration of social information that enforces market dependence by restricting the topological degrees of freedom available within the social plenum. Drawing on Mattei's ``austerity trinity'' and her concept of the Capital Order, the essay constructs a field-theoretic model in which institutional boundaries generate persistent nonlocal signals that continuously deform the scalar potential of the social field---channeling vector flow toward configurations compatible with private ownership and wage dependence, and suppressing high-entropy alternatives. The Spherepop calculus of events formalizes the irreversibility of austerity measures, while the RSVP potential decomposition provides a structural account of technocratic persistence. The concluding section argues that economic freedom, understood as re-democratization, corresponds formally to the restoration of high-entropy potential and the expansion of admissible topological configurations within the plenum.
\end{abstract}

\newpage

\section{Introduction: The Political Calculus of Economic Constraints}

In \emph{The Capital Order} (2022) and its companion volume \emph{Escape from Capitalism}, Clara E.\ Mattei argues that austerity is not an economic inevitability but a technology of class discipline: a ``vital bulwark in defense of the capitalist system'' deployed by economic experts to silence calls for democratic alternatives and reimpose structural dependence on wage labor. Her historical analysis centers on the post-WWI moment when the perceived naturalness of capitalism fractured---when wartime collectivization revealed that private ownership of production was a political choice rather than an economic law---and when a coalition of technocrats in Britain and Italy systematically invented what she terms the ``austerity trinity'' to restore the conditions of that choice.

This essay takes Mattei's framework as its political-economic starting point and asks a structural question: what kind of formal system does the Capital Order constitute? The thesis advanced here is that austerity, understood as a comprehensive project of constraint, can be rigorously modeled as a low-entropy attractor configuration in the Relativistic Scalar-Vector Plenum (RSVP) framework---a field-theoretic system defined by a scalar density field $\Phi$, a vector flow field $\mathbf{v}$, and an entropy field $S$---with the irreversibility of its measures formalized through the Spherepop calculus of append-only events. The goal is not to reduce political economy to physics, but to use the precision of field theory and irreversible computation to expose the structural geometry of how austerity operates: not as an external imposition but as a deformation of the very landscape within which social possibility unfolds.

\section{RSVP and the Field Theory of Class Discipline}

\subsection{The Scalar Potential of Market Dependence}

Within the RSVP framework, the social plenum is modeled as a continuous field over a domain $\Omega \subset \mathbb{R}^d$, where the scalar potential $\Phi(x,t)$ encodes the accessibility of resources and decision-making power at each point in the field, the vector field $\mathbf{v}(x,t)$ represents the directed flow of human effort, attention, and economic agency, and the entropy field $S(x,t)$ measures the degree of available variation---the effective dimensionality of social possibility---at each location.

Mattei's central claim is that the Capital Order is built on two mutually reinforcing pillars: private ownership, which concentrates decision-making power in the hands of those who control the means of production, and market dependence, which compels the majority to sell their labor for a wage because no alternative access to subsistence exists. In RSVP terms, this corresponds to a configuration in which the scalar potential $\Phi$ exhibits a steep gradient directed toward capital accumulation, while the entropy field $S$ is suppressed in the regions of the plenum occupied by the working class.

Austerity, in this mapping, functions by lowering what we can call the \emph{subsistence potential}---the minimum threshold of $\Phi$ at which autonomous economic action becomes possible. When public expenditure on health, education, and welfare is cut (fiscal austerity), the scalar potential floor is depressed across large regions of the plenum, steepening the gradient that drives vector flow toward wage labor. When interest rates are raised and the money supply contracted (monetary austerity), the dissipative coefficient $\nu$ effectively increases for workers while creditors are insulated from entropy. When wages are suppressed and unions dismantled (industrial austerity), the boundary conditions on $\mathbf{v}$ are tightened, restricting the admissible directions of collective agency.

\subsection{Entropy Fields and the Suppression of Alternatives}

Mattei's description of the post-WWI ``crisis of capitalism'' maps precisely onto a transition in the entropy field. Wartime collectivization---the state management of mines, railways, and food production---demonstrated to working-class movements across Europe that the organization of production was not governed by natural law but by political choice. In RSVP terms, this was a period of elevated $S$: the effective dimensionality of the social possibility space expanded as the ``numbing story'' of capitalist inevitability broke down, and movements such as Italy's \emph{L'Ordine Nuovo} and Britain's shop steward councils began to populate regions of the plenum previously considered inaccessible.

Austerity, in this context, functions as an \emph{entropic dampener}. By deploying technocratic authority---the insulation of economic decision-making from democratic participation through institutions like central banks and international credit bodies---the experts Mattei documents worked to collapse the expanded possibility space back into a singular, low-entropy attractor: the configuration of wage labor under private ownership. The ``pure economics'' they employed was not a neutral science but a representational technology that framed this attractor as the unique equilibrium of a natural system, thereby naturalizing what was in fact a political stabilization of the field.

Formally, the dampening operation can be represented as a reduction in the entropy gradient $\nabla S$ available to the working-class regions of $\Omega$, combined with an increase in the curvature of $\Phi$ around the wage-labor attractor---making local departures from that configuration energetically costly and globally unstable. Unemployment, in this model, is not a side effect of austerity but an instrument: it raises the effective ``restoring force'' that returns displaced workers to the wage-labor configuration, functioning precisely as a high-gradient region in $\Phi$ that channels flow back toward market dependence.

\section{Spherepop: Computation as Social Stratification}

\subsection{The Irreversibility of Austerity Measures}

The Spherepop framework models computation as an irreversible calculus of events: each committed event appends to an immutable history and creates a new topological scope---a ``bubble''---that restricts the admissible operations available within its domain. Unlike reversible computation, where states can be freely traversed in either direction, Spherepop captures systems in which the past is structurally determinative of the present: prior events narrow the grammar of what can be said or done next.

This architecture provides an unusually precise formalization of how austerity operates historically. When a social service is cut, a public enterprise is privatized, or a union's legal standing is revoked, the event is not merely a policy change that could be symmetrically reversed. It commits a new topological configuration to the plenum: workers who depended on the service now face restructured dependencies; the institutional knowledge and collective organization embedded in the public enterprise is dissolved; the legal grammar governing labor relations is rewritten with the union's capacity removed. Each such event creates a bubble within which future political and economic action must occur, and the grammar of that bubble is defined by the committed configuration.

Mattei documents this process with particular clarity in her account of the Italian stabilization of 1926 and the British return to gold in 1925. In both cases, the decisive moves were not merely economic but constitutional: they embedded austerity in institutional structures---the independence of the central bank, the convertibility regime, the criminalization of strikes---that placed economic governance beyond the reach of ordinary democratic contestation. In Spherepop terms, these were scope-defining events that restructured the topological space within which all subsequent social action would unfold.

\subsection{Nested Scopes and the Technocratic Shield}

Mattei argues that a defining feature of modern austerity is the delegation of economic authority to technocratic institutions whose internal grammar is insulated from the democratic plenum. Central banks, international financial institutions, and bodies of credentialed economic experts constitute, in her analysis, a layer of decision-making that operates above and prior to the political choices available to citizens.

In the Spherepop architecture, this corresponds to a structure of nested scopes in which the innermost and most consequential bubbles are populated by actors whose operations are invisible to---and cannot be overwritten by---the outer scopes in which ordinary political life takes place. Democratic elections, legislative debates, and social movements all occur within bubbles that are themselves nested inside the technocratic scope. The grammar of that outer scope---the assumption that inflation must be controlled, that sovereign debt must be serviced, that the money supply is a technical rather than political variable---is not available for revision from within the inner scopes.

This is precisely what Mattei means when she argues that austerity ``excludes the public from economic decision-making'' and delegates power to institutions whose authority rests on the appearance of scientific neutrality. The ideology of ``pure economics'' is, in Spherepop terms, the grammar of the outer scope presented as if it were the structure of reality itself---as if the rules of the bubble were the rules of the field.

\section{The Nonlocal RSVP Model of Institutional Forcing}

\subsection{The Potential Landscape of Modern Austerity}

Within the RSVP framework, the dynamics of the social plenum are governed by the evolution of a vector field $\mathbf{v}(x,t)$ over a scalar potential $\Phi(x,t)$. We consider a domain $\Omega \subset \mathbb{R}^d$ representing the social plenum, with boundary $\partial\Omega$ corresponding to the institutional surface of the Capital Order. The evolution of the vector field is given by

\begin{highlight}
\begin{align*}
\partial_t \mathbf{v} + (\mathbf{v}\cdot\nabla)\mathbf{v}
&\;=\;
-\nabla\Phi + \nu\,\Delta\mathbf{v}, \\[4pt]
\nabla\cdot\mathbf{v} &\;=\; 0,
\end{align*}
\end{highlight}

where $\nu > 0$ encodes dissipative structure. The scalar potential is decomposed as
\[
\Phi(x,t) = \Phi_P(x,t) - \Phi_S(x,t) - \Phi_B(x,t),
\]
where $\Phi_P$ represents pressure-like structural constraints, $\Phi_S$ encodes entropic structure, and $\Phi_B$ is a nonlocal boundary-driven contribution.

The entropic component is modeled as
\[
\Phi_S(x,t)
=
\sum_{n=0}^{N}
\frac{E_n(x,t)}{mc^2}\;k\;x_n(x,t)\bigl(1-x_n(x,t)\bigr)\ln P_x(x,t),
\]
which captures bounded mixture dynamics modulated by a logarithmic information term. The logistic factor $x_n(1-x_n)$ reaches its maximum at $x_n = \tfrac{1}{2}$---the point of greatest internal variation---and vanishes at the boundary values $x_n \in \{0,1\}$ corresponding to complete exclusion from, or saturation of, a given resource configuration. Austerity, in this term, operates by driving the $x_n$ toward their boundary values, collapsing the internal variation of the system and suppressing $\Phi_S$.

The boundary contribution is defined through a global functional that aggregates structured activity on the institutional boundary:

\begin{highlight}
\[
Q(t)
=
\operatorname{Re}
\oint_{\partial\Omega}
e^{i(V-E+F)}
\Bigl[
\partial_t \Psi
+
i\Psi(\partial_t V-\partial_t E+\partial_t F)
\Bigr]\,ds,
\]
\end{highlight}

where $\Psi$, $V$, $E$, $F$ are real scalar fields defined on $\partial\Omega$ encoding the institutional configuration---the state of financial regulation, labor law, creditor protection, and technocratic authority at the boundary surface. This signal is injected into the interior via a spatial profile $\phi:\Omega\to\mathbb{R}$,
\[
\Phi_B(x,t) = \alpha\,\phi(x)\,Q(t),
\]
with coupling strength $\alpha \in \mathbb{R}$. Austerity is therefore not a localized intervention but a boundary-driven modulation of the entire potential field. The Capital Order acts as an active informational surface that continuously generates signals constraining the admissible flow of the system.

\subsection{Incompressibility and Class Rigidity}

The incompressibility condition $\nabla\cdot\mathbf{v}=0$ imposes a conservation constraint on the flow of the social plenum. In physical systems, this enforces volume preservation; in the present model, it encodes the structural rigidity of class relations under austerity. Resources and agency that flow away from one region of the plenum must appear in another: the system is closed under redistribution of flow, meaning that gains in one sector are structurally tied to losses in another.

Taking the divergence of the momentum equation yields a Poisson equation for the pressure component,
\[
\Delta\Phi_P
=
-\nabla\cdot\bigl((\mathbf{v}\cdot\nabla)\mathbf{v}\bigr)
+\Delta\Phi_S
+\Delta\Phi_B.
\]
This relation shows that the pressure field is not independently imposed but is determined by the global configuration of the system. Any attempt to locally alter the flow---through redistributive policies, public investment, or democratic economic planning---induces an immediate global adjustment mediated by the elliptic operator. The system responds to local perturbations nonlocally: the correction propagates across the entire domain $\Omega$ instantaneously.

The presence of the boundary term $\Phi_B$ ensures that this global adjustment is not neutral. Instead, the institutional boundary injects corrective signals that restore the system toward configurations compatible with the Capital Order. Class rigidity emerges in this model not as a contingent enforcement of rules but as a structural consequence of nonlocal constraint propagation: the geometry of the system itself resists departures from the low-entropy attractor.

\subsection{Boundary Memory as Technocratic Persistence}

The functional $Q(t)$ defines a form of boundary memory. Because it depends on the temporal derivatives $\partial_t\Psi$, $\partial_t V$, $\partial_t E$, $\partial_t F$---that is, on the rate of change of the institutional configuration---it encodes the accumulated historical momentum of policy, regulation, and technocratic intervention. Rapid institutional change produces large signals; gradual drift produces small ones. But in both cases, past boundary activity is continuously reintroduced into the present dynamics of the system through the coupling $\Phi_B(x,t) = \alpha\,\phi(x)\,Q(t)$.

In the language of Spherepop, each modification to the boundary commits a new event to the append-only history of the plenum. These events do not expire: their effects persist in $Q(t)$ and continue to shape the potential landscape. This provides a structural explanation for what Mattei observes historically---the persistence of austerity even when its stated objectives, such as restoring growth or reducing debt, repeatedly fail. The failure of the stated justification does not diminish the boundary signal. The institutional structures that generate $Q(t)$---central bank independence, debt-servicing obligations, austerity-conditioned international credit---continue to inject restrictive terms into $\Phi_B$ regardless of whether austerity is delivering on its announced economic premises.

This corresponds, in Mattei's analysis, to her observation that austerity functions not as a policy instrument optimized for growth but as a ``bulwark'' whose primary function is to maintain the conditions of the Capital Order. The boundary memory formalizes this: past configurations constrain future possibilities not through ongoing enforcement but through the geometric structure of the field itself. The Capital Order persists because it is embedded in the boundary, and the boundary continuously speaks into the interior.

\section{The Mechanics of Escape: Re-Democratization as Topological Expansion}

Mattei's concept of ``escape'' from the Capital Order is defined not as a single transformative event but as a shift in the epistemic and political conditions under which economic decisions are made---a re-democratization of the economy that restores collective agency over the potential landscape. In RSVP terms, this corresponds to an expansion of the entropy field $S$ across the working-class regions of $\Omega$, combined with a restructuring of the boundary conditions on $\partial\Omega$ such that the signals $Q(t)$ no longer exclusively encode configurations compatible with private ownership and market dependence.

The specific intervention Mattei advocates---using the lens of class to ``re-see'' the world, acknowledging economic decisions as political acts, and reclaiming participatory governance of the potential landscape---maps onto a sequence of changes in the field geometry. First, making $Q(t)$ democratically responsive requires altering the institutional configuration on $\partial\Omega$: the structures that currently insulate technocratic signals from democratic input must be opened to the plenum's interior. Second, expanding $\Phi_S$ requires policies that increase the logistic variation $x_n(1-x_n)$---raising the subsistence potential floor so that workers are not driven to the boundary values of complete dependence. Third, weakening the coupling $\alpha$ between the institutional boundary and the interior---or redistributing it so that $\phi(x)$ amplifies high-entropy rather than low-entropy regions---changes the direction in which boundary memory flows through the system.

In Spherepop terms, re-democratization requires inserting new events into the boundary history that alter the grammar of the outer scope---that make visible, and therefore revisable, the rules that currently present themselves as natural constraints. Each successful assertion of democratic economic governance is, in this model, a scope-redefining event: it does not merely change a policy within the existing structure but modifies the topological space within which future decisions are made.

The formal implication is that the potential landscape of the social plenum is not fixed by nature. It is maintained by the continuous injection of boundary signals into the scalar potential, and it can be restructured by changing what the boundary encodes. Economic freedom, as Mattei defines it---emancipation from exploitation, the capacity to prioritize human needs over the logic of profit---is in RSVP terms the state in which the boundary generates high-entropy signals that amplify rather than suppress the admissible configurations of the plenum.

\section{Conclusion}

This essay has argued that austerity, understood through Mattei's historical and theoretical analysis, constitutes a geometric constraint on the social plenum: a systematic deformation of the scalar potential $\Phi$, a suppression of the entropy field $S$, and a structuring of the institutional boundary $\partial\Omega$ such that nonlocal signals continuously channel the vector flow $\mathbf{v}$ toward configurations compatible with private ownership and wage dependence. The RSVP framework provides the formal language for this analysis, while the Spherepop calculus of irreversible events accounts for the historical persistence of austerity---its capacity to continue shaping the field even when its stated economic rationale has collapsed.

The Capital Order is not a law of nature but a particular low-entropy attractor of a social field whose geometry is maintained by political choices embedded in institutional boundaries. The ``pure economics'' that represents this attractor as natural is, in formal terms, the grammar of the outer Spherepop scope presented as if it were the structure of reality. Economic freedom is therefore not utopian but topological: it is the expansion of the admissible phase space of the social plenum, achieved by restructuring the boundary conditions that currently constrain it.

What a field-theoretic approach adds to Mattei's already rigorous political economy is this: it makes the mechanism of constraint precise. Austerity does not suppress alternatives by forbidding them; it suppresses them by reshaping the potential landscape such that alternatives require increasingly large departures from the attractor to reach. The boundary memory ensures that this shaping is continuous and persistent. And the incompressibility condition ensures that any local gain against the attractor propagates globally, triggering nonlocal corrective responses. This is the geometry of the Capital Order. Making it visible is the first step toward altering it.

\vspace{12pt}

\begin{thebibliography}{99}

\bibitem{mattei2022capital}
Mattei, Clara E.
\textit{The Capital Order: How Economists Invented Austerity and Paved the Way to Fascism}.
Chicago: University of Chicago Press, 2022.

\bibitem{mattei2026escape}
Mattei, Clara E.
\textit{Escape from Capitalism: An Intervention}.
New York: Simon \& Schuster, 2026.

\bibitem{polanyi1944}
Polanyi, Karl.
\textit{The Great Transformation: The Political and Economic Origins of Our Time}.
Boston: Beacon Press, 1944.

\bibitem{keynes1936}
Keynes, John Maynard.
\textit{The General Theory of Employment, Interest, and Money}.
London: Macmillan, 1936.

\bibitem{kalecki1943}
Kalecki, Micha\l.
``Political Aspects of Full Employment.''
\textit{Political Quarterly} 14, no.\ 4 (1943): 322--331.

\bibitem{graeber2011}
Graeber, David.
\textit{Debt: The First 5000 Years}.
Brooklyn: Melville House, 2011.

\bibitem{piketty2014}
Piketty, Thomas.
\textit{Capital in the Twenty-First Century}.
Cambridge, MA: Harvard University Press, 2014.

\bibitem{harvey2005}
Harvey, David.
\textit{A Brief History of Neoliberalism}.
Oxford: Oxford University Press, 2005.

\bibitem{foucault2008}
Foucault, Michel.
\textit{The Birth of Biopolitics: Lectures at the Coll\`ege de France, 1978--1979}.
New York: Palgrave Macmillan, 2008.

\bibitem{streeck2014}
Streeck, Wolfgang.
\textit{Buying Time: The Delayed Crisis of Democratic Capitalism}.
London: Verso, 2014.

\bibitem{mankiw2010}
Mankiw, N.\ Gregory.
\textit{Macroeconomics}.
7th ed.\ New York: Worth Publishers, 2010.

\bibitem{blanchard2017}
Blanchard, Olivier.
\textit{Macroeconomics}.
7th ed.\ Boston: Pearson, 2017.

\bibitem{samuelson1948}
Samuelson, Paul A.
\textit{Economics: An Introductory Analysis}.
New York: McGraw-Hill, 1948.

\bibitem{arrow1951}
Arrow, Kenneth J.
\textit{Social Choice and Individual Values}.
New York: Wiley, 1951.

\bibitem{debreu1959}
Debreu, G\'erard.
\textit{Theory of Value: An Axiomatic Analysis of Economic Equilibrium}.
New Haven: Yale University Press, 1959.

\bibitem{vonNeumann1944}
von Neumann, John, and Oskar Morgenstern.
\textit{Theory of Games and Economic Behavior}.
Princeton: Princeton University Press, 1944.

\bibitem{shannon1948}
Shannon, Claude E.
``A Mathematical Theory of Communication.''
\textit{Bell System Technical Journal} 27 (1948): 379--423, 623--656.

\bibitem{jaynes1957}
Jaynes, Edwin T.
``Information Theory and Statistical Mechanics.''
\textit{Physical Review} 106, no.\ 4 (1957): 620--630.

\bibitem{prigogine1984}
Prigogine, Ilya, and Isabelle Stengers.
\textit{Order Out of Chaos: Man's New Dialogue with Nature}.
New York: Bantam Books, 1984.

\bibitem{haken1983}
Haken, Hermann.
\textit{Synergetics: An Introduction}.
Berlin: Springer, 1983.

\bibitem{anderson1972}
Anderson, Philip W.
``More Is Different.''
\textit{Science} 177, no.\ 4047 (1972): 393--396.

\bibitem{smolin2006}
Smolin, Lee.
\textit{The Trouble with Physics}.
Boston: Houghton Mifflin, 2006.

\bibitem{jacobson1995}
Jacobson, Ted.
``Thermodynamics of Spacetime: The Einstein Equation of State.''
\textit{Physical Review Letters} 75, no.\ 7 (1995): 1260--1263.

\bibitem{verlinde2011}
Verlinde, Erik.
``On the Origin of Gravity and the Laws of Newton.''
\textit{Journal of High Energy Physics} 2011, no.\ 4 (2011): 29.

\end{thebibliography}

\end{document}
